% Generated draft: requires mathematical and source review.
\documentclass[11pt]{article}
\usepackage[margin=25mm]{geometry}
\usepackage{fontspec}
\setmainfont{Latin Modern Roman}
\usepackage{amsmath,amssymb,mathtools,unicode-math}
\setmathfont{Latin Modern Math}
\usepackage{enumitem,needspace,xurl,hyperref}
\setlength{\parindent}{0pt}
\setlength{\parskip}{5pt}
\setlength{\emergencystretch}{3em}
\newcommand{\N}{\mathbb{N}}
\newcommand{\Z}{\mathbb{Z}}
\newcommand{\Q}{\mathbb{Q}}
\newcommand{\R}{\mathbb{R}}
\newcommand{\C}{\mathbb{C}}
\newcommand{\F}{\mathbb{F}}
\newcommand{\A}{\mathbb{A}}
\newcommand{\PP}{\mathbb{P}}
\newcommand{\E}{\mathbb{E}}
\newcommand{\eps}{\varepsilon}
\newcommand{\dd}{\,\mathrm{d}}
\begin{document}
\section*{UnsolvedMath ID: 5000006}
\textbf{Source problem code:} catalogue entry
\subsection{Definitions and mathematical statement}
Let $n\geq3$ and unfold billiard paths in a regular Euclidean $n$-gon by reflecting the polygon across every encountered side. A short trajectory is a straight segment in the unfolding that joins two polygon vertices and encounters no vertex in its interior. Its length is the Euclidean length of this segment. A reachable point is an endpoint vertex arising in such an unfolding.

The source assigns an oriented short trajectory source and target endpoint angles $\alpha,\beta$ according to its Definition 2.1. These produce $n-2$ types $A_k$, indexed modulo $n-2$; the supplied extract specifies only that type $A_0$ has $\beta-\alpha=2\pi/n$. It does not include Definition 2.1 or the remaining type-assignment rule.

For parallel short trajectories $\gamma,\eta$ of types $A_k,A_\ell$, respectively, prove
\[
\frac{\operatorname{length}(\gamma)}{\operatorname{length}(\eta)}
=\frac{\sin((k+1)\pi/n)}{\sin((\ell+1)\pi/n)}.
\]
The formula is read using fixed representatives $k,\ell\in\{0,\ldots,n-3\}$ of the type indices. Representatives must be fixed: the displayed sine expression is not intrinsically a function of arbitrary integer representatives modulo $n-2$. The conjecture quantifies over every admissible parallel pair in the source's sense, not merely the existence of a pair with the displayed ratio. Exact endpoint-angle and parallelism conventions require the missing source definition.
\subsection{Short English statement}
In unfolded regular-polygon billiards, prove that parallel vertex-to-vertex trajectories of types $A_k$ and $A_\ell$ have lengths in the ratio $\sin((k+1)\pi/n):\sin((\ell+1)\pi/n)$. The extract omits the full type definition.
\subsection{Sources}
\begin{itemize}
\item Dmitry Fuchs, \emph{Billiard Trajectories in Regular Polygons and Foliations by Closed Geodesics on Surfaces} (2020), Conjecture 2.7 and the referenced Definition 2.1, \url{https://amj.math.stonybrook.edu/pdf-Springer-final/020-0170.pdf}; accessed 2026-07-29, with PDF-text extraction recorded in the supplied provenance.
\item Local registry \texttt{data/unsolvedmath/problems.json}, permanent UnsolvedMath ID $5000006$, source code catalogue entry, \emph{Length ratios of parallel short trajectories}. Dataset attribution: Ulam AI and the UnsolvedMath Contributors.
\end{itemize}
\end{document}
